% P16: source-specific mediation and a measurable-kernel completion.
\section{Causal mediation: direct, indirect and interventional effects}
\label{sec:p16-causal-mediation}

Pearl's original UAI treatment distinguishes controlled and natural effects,
and gives experimental identification before observational identification
\cite[Secs.~3.1--3.7]{Pearl2001DirectIndirectEffects}.
Its notation \(X,Z,W\) corresponds here to exposure \(A\), selected mediator
\(M\), and baseline covariates \(C\). In its Markovian corollary, \(Z\)
contains \emph{all} parents of the response other than \(X\). Effects
relative to a smaller mediator set need not be effects of a single direct
edge. The distinction matters when exposure changes another common cause
of mediator and response.

The interventional single-mediator continuation uses independent mediator
draws and is attributed separately
\cite[author manuscript pp.~2--4]{VansteelandtDaniel2017InterventionalEffects}.
All effects below are population functionals, with additive contrasts.
No estimator fitted to observations is identified with its population target.

\subsection{Jointly defined nested counterfactuals}

\begin{hbdefinition}{def:p16-mediation-worlds}{A mediator and a jointly measurable response family}
In the finite structural setting of \cref{def:p16-counterfactual-worlds},
take disjoint exposure, mediator and response coordinates and baseline
\(C\) unaffected by exposure or mediator interventions. Fix two distinct
exposure values \(a_0,a_1\). Write
\[
 M_a=M[A=a],\qquad Y_{a,m}=Y[A=a,M=m],\qquad
 Y_{a,M_b}(\omega)=Y_{a,m}(\omega)\big|_{m=M_b(\omega)}.
\]
The last expression evaluates a finite table on the same outcome; it is
not a new assumption of independence between worlds.

For general spaces, require \(C,A,M,R\) to take values in nonempty
standard Borel spaces, where standard Borel means admitting a measurable
isomorphism to a Borel subset of a complete separable metric space.
Require jointly measurable maps \(m(a,\omega)\) and
\(f(a,m,c,r)\), with \(f\) real-valued, on one probability space
\((\Omega,\mathcal F,P)\).
Define \(M_a=m(a,\omega)\),
\(Y_{a,m}=f(a,m,C,R)\), and the nested variable by measurable composition
\[
 Y_{a,M_b}=f(a,m(b,\omega),C(\omega),R(\omega)).
\]
The random element \(R\) is a specified response background; it may
encode more information than any one potential response. Assume pointwise
consistency \(M=M_A\), \(Y=Y_{A,M}\), and composition
\(Y_a=Y_{a,M_a}\). The baseline is the same coordinate in all worlds.
For mean effects assume the relevant absolute expectations are finite.
Statements about distributions instead apply to bounded measurable
response functions, without an integrability assumption on \(Y\).
\end{hbdefinition}

\begin{hbtheorem}{thm:p16-mediation-measurability}{proved=annex:p16.mediation-measurability}{Structural meaning and measurable random substitution}
The finite nested variables are well defined and composition follows from
\cref{thm:p16-counterfactual-laws}. More generally a finite acyclic graph
with standard Borel coordinate/background spaces and jointly measurable
total structural functions has jointly measurable solutions in its
intervention values and background. It supplies the displayed nested
variables, consistency and composition. No equality of conditional versions
at separate fixed mediator values is used in this construction.
\end{hbtheorem}

\begin{hbdefinition}{def:p16-mediation-effects}{Controlled and four natural contrasts}
Put \(\theta_{ij}=\mathbb E[Y_{a_i,M_{a_j}}]\), for \(i,j\in\{0,1\}\).
The controlled direct effect at a prescribed mediator value is
\(\operatorname{CDE}(m)=\mathbb E[Y_{a_1,m}-Y_{a_0,m}]\).
The unit effect is the same difference before expectation; qualitative
presence means that difference is nonzero. Define
\[
\begin{array}{ll}
 \operatorname{TE}=\theta_{11}-\theta_{00},&
 \operatorname{PNDE}=\theta_{10}-\theta_{00},\\
 \operatorname{TNDE}=\theta_{11}-\theta_{01},&
 \operatorname{PNIE}=\theta_{01}-\theta_{00},\\
 \operatorname{TNIE}=\theta_{11}-\theta_{10}.&
\end{array}
\]
Here P and T distinguish the pure and total versions. Pearl's forward
\(\operatorname{NDE}(a_1,a_0)\) is PNDE; his forward
\(\operatorname{NIE}(a_1,a_0)\) is PNIE. The later continuation's indirect
effect paired with PNDE is TNIE. These naming conventions are not interchangeable.
For unit contrasts replace each \(\theta_{ij}\) by the corresponding
\(Y_{a_i,M_{a_j}}(\omega)\). Qualitative natural direct or indirect
presence means the corresponding pure unit contrast is nonzero. A
controlled direct effect exists at a unit if some prescribed mediator
value gives a nonzero controlled contrast.
\index{natural direct effect}\index{natural indirect effect}
\end{hbdefinition}

\begin{hbtheorem}{thm:p16-mediation-decompositions}{proved=annex:p16.mediation-decompositions}{The two decompositions and the interaction remainder}
Pointwise and in expectation,
\[
 \operatorname{TE}=\operatorname{PNDE}+\operatorname{TNIE}
                  =\operatorname{TNDE}+\operatorname{PNIE}.
\]
If \(I=\theta_{11}-\theta_{10}-\theta_{01}+\theta_{00}\), then
\[
 \operatorname{TNDE}-\operatorname{PNDE}
 =\operatorname{TNIE}-\operatorname{PNIE}=I,
 \qquad\operatorname{TE}=\operatorname{PNDE}+\operatorname{PNIE}+I.
\]
Equivalently Pearl's original reverse-transition identities are
\(\operatorname{TE}(a_1,a_0)=\operatorname{NIE}(a_1,a_0)
-\operatorname{NDE}(a_0,a_1)\) and
\(\operatorname{TE}(a_1,a_0)=\operatorname{NDE}(a_1,a_0)
-\operatorname{NIE}(a_0,a_1)\).
If \(Y_{a,m}=u(a,C,R)+v(m,C,R)\), then \(I=0\).
This includes additive linear structural responses. In general the two
forward pure effects alone do not add to the total effect.
\end{hbtheorem}

\subsection{Experimental identification and the general-space qualification}

\begin{hbtheorem}{thm:p16-mediation-countable-formula}{proved=annex:p16.mediation-countable-formula}{Pearl's experimental factorization for a countable mediator}
Suppose the mediator value set is finite or countable. For each requested
pair \((a,b)\), assume
\(Y_{a,m}\mathrel{\perp\!\!\!\perp}M_b\mid C\) for every mediator
value \(m\). Let \(\eta_b(c,\cdot)\) be a conditional law of \(M_b\)
given \(C=c\), and \(Q_a(m,c,\cdot)\) a conditional law of \(Y_{a,m}\)
given \(C=c\). Then for every bounded measurable \(\varphi\),
\[
 \mathbb E[\varphi(Y_{a,M_b})]
 =\int\sum_m\left\{\int\varphi(y)Q_a(m,c,dy)\right\}
                  \eta_b(c,\{m\})\,\nu(dc),                 \tag{MF}
\]
where \(\nu\) is the baseline law. For means require the corresponding
absolute integral to be finite. Knowledge of these intervention kernels
on their weighted support identifies the nested mean and distribution.
Their identification may come from separate randomized interventions or
from any separately justified observational identification argument.

For both \(a=a_1,a_0\) with \(b=a_0\), the formula yields Pearl's
\[
 \operatorname{PNDE}=\int\sum_m
 [q_{a_1}(m,c)-q_{a_0}(m,c)]\eta_{a_0}(c,\{m\})\nu(dc),
\]
where \(q_a(m,c)=\int yQ_a(m,c,dy)\) on the measurable set with
finite absolute integral, and is assigned zero elsewhere. The analogous PNIE formula
uses \(q_{a_0}\) and the signed difference \(\eta_{a_1}-\eta_{a_0}\).
The assumptions must cover every factor actually used in that subtraction.
\end{hbtheorem}

\begin{hbdefinition}{def:p16-mediation-uniform-kernels}{One conditional product law for the response background}
For general standard Borel mediator spaces let \(\kappa(c,dr)\) be a
conditional law of \(R\) given \(C\). A sufficient \emph{uniform
cross-world condition} for the pair \((a,b)\) is
\[
 \mathcal L(M_b,R\mid C=c)=\eta_b(c,dm)\kappa(c,dr)          \tag{CW}
\]
outside one \(\nu\)-null set, as equality of probability measures on
the product space. Conditional independence of these random elements
has exactly this meaning. It concerns the specified response background,
not separate versions for individual \(m\). Put
\[
 Q_a(m,c,B)=\int\mathbf1_B(f(a,m,c,r))\kappa(c,dr).
\]
These are one jointly measurable family of probability kernels.
\end{hbdefinition}

\begin{hbtheorem}{thm:p16-mediation-kernel-formula}{proved=annex:p16.mediation-kernel-formula}{The mediation formula with a justified random mediator}
Under (CW), for every bounded measurable \(\varphi\),
\[
 \mathbb E[\varphi(Y_{a,M_b})]
 =\int\!\int\!\int\varphi(y)Q_a(m,c,dy)\eta_b(c,dm)\nu(dc).
\]
The inner mediator mixture is a conditional probability kernel for the
nested response given \(C\), with one common null set for equality as
probability measures. For the identity function require
\(\int\!\int\!\int|f(a,m,c,r)|\kappa(c,dr)\eta_b(c,dm)\nu(dc)<\infty\).
The product condition is a sufficient general-space completion of the
source's discrete summation proof. Pointwise conditional independence
\(Y_{a,m}\perp M_b\mid C\) at every fixed \(m\), even with a jointly
measurable response family, does not imply this theorem for an uncountable
mediator. \Cref{thm:p16-mediation-diagonal} proves the failure explicitly.
\end{hbtheorem}

\subsection{Observational identification and support}

\begin{hbtheorem}{thm:p16-mediation-finite-observed}{proved=annex:p16.mediation-finite-observed}{The finite observational mediation formula}
In the finite model assume, for all values needed in the target,
\[
\begin{aligned}
 Y_{a,m}&\mathrel{\perp\!\!\!\perp}A\mid C,&
 M_b&\mathrel{\perp\!\!\!\perp}A\mid C,\\
 Y_{a,m}&\mathrel{\perp\!\!\!\perp}M\mid(A=a,C),&
 Y_{a,m}&\mathrel{\perp\!\!\!\perp}M_b\mid C.
\end{aligned}
\]
For every \(c\) of positive baseline mass require
\(p(A=a\mid c)>0\), \(p(A=b\mid c)>0\), and
\[
 p(M=m\mid A=b,c)>0\Longrightarrow p(M=m\mid A=a,c)>0.
\]
Then
\[
 \mathbb E[Y_{a,M_b}]
 =\sum_{c,m} \mathbb E[Y\mid A=a,M=m,C=c]
                       p(M=m\mid A=b,c)p(c),               \tag{OMF}
\]
omitting zero-weight cells before requesting conditional means. Without
the last, cross-world independence condition, the same expression can
identify an independent mediator-draw target instead. It does not thereby
identify the natural effect. A controlled mean at \(m\) uses the first
and third independence conditions and requires that mediator value to
have positive probability at every contributing \((a,c)\).
\end{hbtheorem}

\begin{hbdefinition}{def:p16-mediation-observed-general}{Uniform observational conditions with an intermediate confounder}
Permit a recorded variable \(L\) between exposure and mediator. Specify
a jointly measurable map \(\ell(a,\omega)\), put \(L_a=\ell(a,\omega)\),
and require pointwise \(L=L_A\) and no mediator effect on \(L\). All coordinate spaces
are standard Borel and all relevant maps are jointly measurable.
For each target exposure value assume positive conditional atom
\(e_a(c)=\mathbb P(A=a\mid C=c)>0\) for \(\nu\)-almost every \(c\).
Let \(\lambda_a(c,dl)=\mathcal L(L\mid A=a,C=c)\),
\(\rho_a(c,l,dm)=\mathcal L(M\mid A=a,C=c,L=l)\), and let
\(\mu_a(c,l,m)\) be the mean of one regular conditional response kernel
given \((A=a,C=c,L=l,M=m)\), assigned zero on the measurable set where
its absolute integral is infinite.

Assume \(R\perp A\mid C\), \(M_b\perp A\mid C\), and
\(R\perp M\mid(A=a,C,L)\), as conditional product-kernel identities.
The first and third are uniform response-background exchangeability;
they are stronger than separate conditional assertions for each fixed
mediator. Consistency remains \(Y=f(A,M,C,R)\).
For a known mediator kernel \(k(c,dm)\), require
\[
 k(c,\cdot)\ll\rho_a(c,l,\cdot)
 \quad\hbox{for }\nu(dc)\lambda_a(c,dl)\hbox{-almost every }(c,l). \tag{OV}
\]
Absolute continuity means every set of zero mass for the right-hand
measure has zero mass for the left-hand measure. For a natural or
interventional reference law take \(k=\eta_b\); mediator exchangeability
then identifies \(\eta_b(c,\cdot)=\mathcal L(M\mid A=b,C=c)\), provided
\(e_b(c)>0\). An arbitrary known \(k\) needs no assumption about \(M_b\).
\end{hbdefinition}

\begin{hbdefinition}{def:p16-mediation-policy}{The interventional analogue}
Given a probability kernel \(k\) from baseline space to mediator space,
enlarge the original probability space to have joint law
\(P(d\omega)k(C(\omega),dg)\). Define the new mediator coordinate
\(G^k(\omega,g)=g\) and response \(Y_{a,G^k}=f(a,g,C,R)\).
This retains the original unit distribution and draws the mediator
independently of its remaining background conditional on baseline.
For \(k=\eta_{a_j}\), write \(J_{ij}=\mathbb E[Y_{a_i,G^k}]\).
The interventional direct and indirect contrasts paired in the continuation
are \(J_{10}-J_{00}\) and \(J_{11}-J_{10}\), respectively.
\end{hbdefinition}

\begin{hbtheorem}{thm:p16-mediation-general-observed}{proved=annex:p16.mediation-general-observed}{A common observed kernel and transferred support}
Assume the general observational conditions, \(\mathbb E|Y|<\infty\), and
\(\int\!\int\!\int |f(a,m,c,r)|\kappa(c,dr)k(c,dm)\nu(dc)<\infty\).
Then
\[
 J(a,k)=\int\!\int\!\int
                  \mu_a(c,l,m)k(c,dm)\lambda_a(c,dl)\nu(dc) \tag{IGF}
\]
is the mean response under exposure \(a\) and a mediator drawn from
\(k(C,\cdot)\) independently of the original unit given \(C\).
No natural cross-world independence is needed. With (CW) and
\(k=\eta_b\), it also equals \(\mathbb E[Y_{a,M_b}]\).
For constant \(L\), this is the standard mediation integral replacing
the finite sums in (OMF). For a controlled mediator use \(k=\delta_m\)
when (OV) holds. For non-atomic exposure distributions the conditional-atom
hypothesis does not hold at a prescribed point; fixed-exposure observational
identification then requires additional model or version restrictions.
Experimental intervention kernels in the preceding theorem do not require
an observational exposure atom.
For bounded measurable \(\varphi\), the same identity holds with
\(\varphi(Y_{a,G^k})\) as response and the observed conditional
\(\varphi\)-mean as integrand; the two absolute-integrability assumptions
are then unnecessary.
\end{hbtheorem}

\begin{hbtheorem}{thm:p16-mediation-four-set-criterion}{proved=annex:p16.mediation-four-set-criterion}{The source's four-set graphical criterion for intermediate vertices}
Let \(G\) be a finite acyclic structural graph with independent
disturbances, including any unobserved endogenous common causes.
Let \(A,Y\) be distinct observed vertices and \(M\) a nonempty set of
other observed intermediate vertices, each a descendant of \(A\) and an ancestor
of \(Y\). Mediator values are indexed tuples. Choose observed covariate sets
\(W_0,W_1,W_2,W_3\), excluding \(A,Y\) and every member of \(M\). Write
\(G_{\underline B}\) for deletion of arrows leaving \(B\) and
\(G_{\overline B}\) for deletion of arrows entering \(B\).
Assume the four d-separations
\[
\begin{aligned}
 Y&\perp M\mid W_0 &&\text{in }G_{\underline A\underline M},\\
 Y&\perp A\mid W_0,W_1 &&\text{in }G_{\underline A\overline M},\\
 Y&\perp M\mid A,W_0,W_1,W_2 &&\text{in }G_{\underline M},\\
 M&\perp A\mid W_0,W_3 &&\text{in }G_{\underline A}.
\end{aligned}
\]
Require \(W_0,W_1,W_3\) to contain no descendants of \(A\), and
\(W_2\) no descendants of \(M\). Set \(C=W_0\). Then the
cross-world condition of the finite mediation formula holds, and its
controlled-response and reference-mediator factors are
\[
\begin{aligned}
 q_a(m,c)&=\sum_{w_1,w_2}\mathbb E[Y\mid a,m,c,w_1,w_2]
             p(w_2\mid a,c,w_1)p(w_1\mid c),\\
 \eta_b(c,m)&=\sum_{w_3}p(m\mid b,c,w_3)p(w_3\mid c).
\end{aligned}
\]
Require positive exposure probabilities for \(a\) in every contributing
\((c,w_1)\) stratum and for \(b\) in every contributing \((c,w_3)\)
stratum. Require positive \(p(m\mid a,c,w_1,w_2)\) whenever the
reference mediator cell and both displayed covariate weights are positive.
Zero-weight terms are omitted first. These formulas, inserted into (MF),
identify the natural effect. This theorem states and proves the original
graphical corollary's finite intermediate-variable case separately from the
general-space response-kernel theorem.
\end{hbtheorem}

For the corresponding finite graph with standard Borel coordinate spaces,
\cref{thm:p16-measurable-global-markov} supplies the general graphical
premise and \cref{thm:p16-measurable-mediation-fields,thm:p16-measurable-mediation-corollary} supply the complete uniform
response-field construction and identifying integral. Their conditional
exposure-atom and mediator-support requirements remain explicit.

\begin{hbtheorem}{thm:p16-mediation-parent-scope}{proved=annex:p16.mediation-parent-scope}{The source's independent-error full-parent route}
In a finite acyclic structural model with independent disturbances, let
\(A\) be a parent of \(Y\), and let the mediator tuple contain all
other parents of \(Y\).
For the observational formula require \(A,M,Y\) to be observed.
Then \(Y_{a,m}\) is a function of \(a,m,U_Y\) alone, and is independent
of \(M_b\). On positive observed parent strata its law is
\(\mathcal L(Y\mid A=a,M=m)\). Thus, if a baseline set \(S\) identifies
the mediator law by back-door adjustment,
\[
 \mathbb E[Y_{a,M_b}]
 =\sum_{s,m}\mathbb E[Y\mid a,m]p(m\mid b,s)p(s),
\]
with positivity for every contributing mediator cell and adjustment
stratum. If \(A\) has no confounding for its effect on this tuple, the
mediator factor reduces to \(p(m\mid b)\). This is the source's
Corollary~2 and its simple-model specialization at their full-parent scope.
For standard Borel coordinates the same response-background proof supplies
(CW) and \(Q_a\); a mediator identification formula must separately hold
as an equality of kernels on the target support.
\end{hbtheorem}

\subsection{Independent mediator draws and their decomposition}

\begin{hbtheorem}{thm:p16-mediation-policy-identities}{proved=annex:p16.mediation-policy-identities}{Which total effect the interventional contrasts decompose}
The enlarged measure and response above are well defined. Their mean is
\(\int q_a(m,c)k(c,dm)\nu(dc)\), with no (CW) premise. The two
interventional contrasts sum to \(J_{11}-J_{00}\); the alternative pair
\((J_{11}-J_{01})+(J_{01}-J_{00})\) has that same sum. In general this
is not \(\mathbb E[Y_{a_1}]-\mathbb E[Y_{a_0}]\).
Under (CW) for all four needed pairs, \(J_{ij}=\theta_{ij}\), so all the
corresponding natural and interventional contrasts coincide. If \(L\)
is constant, the finite observational conditions except cross-world
independence hold, and \(Y_{a_i}\perp A\mid C\) for \(i=0,1\),
the diagonal interventional means still equal the
natural single-world means by observed total expectation and exposure
exchangeability. That limited diagonal equality needs no off-diagonal
cross-world assertion.
\end{hbtheorem}

\subsection{Worked effects and failures}

\begin{hbtheorem}{thm:p16-mediation-nonlinear-example}{proved=annex:p16.mediation-nonlinear-example}{An interaction changes both decompositions}
Let \(A\) be fair, \(E\) an independent Bernoulli\((1/4)\) bit,
\(M=A\oplus E\), and \(Y=A+2M+3AM\). The four observed
\((A,M,Y)\) cells are \((0,0,0),(0,1,2),(1,0,1),(1,1,6)\),
with masses \(3/8,1/8,1/8,3/8\). The complete mean table is
\[
 \begin{array}{c|cc}
 \theta_{ij}&j=0&j=1\\\hline
 i=0&1/2&3/2\\i=1&9/4&19/4
 \end{array}
\]
Thus TE is \(17/4\), PNDE is \(7/4\), PNIE is \(1\), TNDE is
\(13/4\), TNIE is \(5/2\), and the interaction remainder is \(3/2\).
The controlled direct effect is \(1+3m\). The mediation formula and
the direct structural calculation agree in every cell.
\end{hbtheorem}

\begin{hbtheorem}{thm:p16-mediation-crossworld-countermodels}{proved=annex:p16.mediation-crossworld-countermodels}{Randomized exposure and all single-world experiments can agree}
There are two finite acyclic models on observed \((A,M,Y)\), with fair
randomized exposure, identical strictly positive observed laws and identical
laws under every intervention on observed coordinates, but PNDE equal
to \(1/2\) in one and zero in the other. Both satisfy the finite
exposure and observed mediator exchangeability conditions, and full
mediator overlap. Their common (OMF) value for PNDE is \(1/4\),
which is their identified interventional direct effect. The natural
cross-world condition fails. Full tables and equations appear in the proof.
\end{hbtheorem}

\begin{hbtheorem}{thm:p16-mediation-induced-confounding}{proved=annex:p16.mediation-induced-confounding}{An exposure-induced mediator--outcome confounder}
Two independent-error models on the fully observed DAG
\(A\to L\to M\to Y\), with \(L\to Y\), have the same strictly
positive observed law and all the same observed-coordinate intervention
laws. Their mediator-relative PNDEs are respectively \(-1/32\) and
\(1/32\), although their total effect is \(3/32\).
Both have \(p(L=1\mid A=0)=1/4\), \(p(L=1\mid A=1)=1/2\),
\(M=L\oplus E_M\), and \(Y=(LM)\oplus E_Y\), where the two noises
are independent Bernoulli\((1/4)\) bits independent of the treatment
and the disturbance assigning \(L\).

The unadjusted mediator formula gives PNDE \(3/64\), which is neither
natural value. The correctly identified independent-draw means are
\[
 \begin{array}{c|cc}
 J_{ij}&j=0&j=1\\\hline
 i=0&19/64&20/64\\i=1&22/64&24/64
 \end{array}
\]
Their overall contrast is \(5/64\), whereas TE is \(6/64\).
Taking \((L,M)\) together as the mediator changes the estimand and gives
natural direct effect zero. The source's full-parent corollary therefore
does not justify discarding \(L\) from its mediator tuple.
\end{hbtheorem}

\begin{hbtheorem}{thm:p16-mediation-diagonal}{proved=annex:p16.mediation-diagonal}{Separate fixed-mediator null sets do not permit random substitution}
There are jointly measurable standard Borel structural models with
randomized binary exposure and a uniform continuous mediator in either
exposure group, satisfying all four pointwise independence conditions of
the finite theorem for every fixed mediator value, that have identical
observed and controlled-intervention laws but different natural direct
effects. One has PNDE one and the other PNDE zero. Thus neither a
measurable response family nor mediator absolute continuity repairs the
uncountable diagonal problem. The uniform kernel assumptions above
explicitly exclude this counterexample.
\end{hbtheorem}

\begin{hbdefinition}{def:p16-mediation-path-policy}{The original edge-subgraph intervention}
For a finite acyclic structural model, a selected edge subgraph \(g\),
and reference exposure \(a_0\), first compute the complete reference
solution \(v^0(\omega)=V[A=a_0](\omega)\). In every structural equation,
retain the current parent input along edges in \(g\), and replace every
other parent input by its component of \(v^0(\omega)\). Denote this
modified family of equations by \(M_g^0\). Its source path-specific
effect is the response contrast in this family between \(A=a_1\)
and \(A=a_0\). The reference inputs are constants within a unit but
random across units; they keep their original shared background.
\end{hbdefinition}

\begin{hbtheorem}{thm:p16-mediation-path-semantics}{proved=annex:p16.mediation-path-semantics}{The source's path construction and its failure example}
The edge-subgraph construction has a unique solution and, in the
standard Borel formulation, a jointly measurable solution. Under
\(A=a_0\) it returns the original reference response. In Pearl's
Figure~3, with \(M=f_M(A,U_M)\), \(L=f_L(M,A,U_L)\),
\(Y=f_Y(M,L,U_Y)\), selecting \(A\to M\to L\to Y\) gives
\[
 f_Y\bigl(M_{a_0},f_L(M_{a_1},a_0,U_L),U_Y\bigr)-Y_{a_0}.
\]
This effect need not be identified even in finite independent-error models.
The preceding induced-confounding pair, with its intermediate-variable
names exchanged and an unused permitted exposure edge added, supplies
such a pair of witnesses. General path-specific identification criteria
are a separate result from these semantics and countermodels.
\end{hbtheorem}
